\sloppy
\chapter{Conclusion and Future Direction}

\section{Conclusion}

This thesis studied a transformer-based approach to fake news detection and compared three models: a purely sequential deep learning model (BiLSTM), a purely transformer-based model (BERT), and a hybrid of the two (BERT + BiLSTM). Chapter~2 reviewed nineteen related works on classical ML, hybrid deep learning, evidence-aware detection, benchmark datasets, and large language models. Chapter~3 proposed a layered, end-to-end detection pipeline with a detailed BERT + BiLSTM classification architecture. Chapter~4 evaluated all three models on the LIAR and ISOT benchmark datasets.

The results are clear. The standalone BiLSTM was inefficient, with only 50\% accuracy and an F1-score of 0.3333. The likely causes are the small number of training epochs, insufficient hyperparameter tuning, and the difficulty of the task. BERT clearly outperformed it, reaching 95\% accuracy and an F1-score of 0.9499, thanks to its pre-trained bidirectional contextual representations. The hybrid BERT + BiLSTM was designed to join BERT's contextual understanding with the BiLSTM's explicit sequence modelling. It matched BERT almost exactly but did not exceed it under the training conditions used here. In this setup the extra sequential component adds no measurable value, and BERT alone can be used for fake news classification without the added complexity of a hybrid.

Within each dataset, the transformer-based models give stable and reliable results across the benchmark datasets, which makes them good candidates for practical use in social media monitoring and automated fact-checking.

A cross-dataset test with cleaned and deduplicated headlines shows the limit of these in-domain scores. Trained on ISOT, BERT reaches 0.984 accuracy on ISOT but only 0.504 on LIAR2, and the GRU behaves the same way (0.532). The models therefore learn dataset-specific patterns more than transferable cues of fake news. Training on ISOT and LIAR2 together keeps ISOT accuracy at 0.984 and gives a LIAR2 AUC of 0.772, but this has not yet been tested on a third dataset. The main work that remains is to close this gap with stronger encoders, article-body input, retrieval-augmented classification, and explainability analysis.

\section{Future Research Direction}

Large Language Models (LLMs) and Explainable AI (XAI) were not part of the system design evaluated in this thesis, but both are clear and promising directions for future work.

Future work should bring more advanced large language models, such as GPT-4, LLaMA~2, Mistral, and Gemini, into the detection pipeline. Retrieval-Augmented Generation (RAG) in particular could let the model query verified databases such as Snopes, PolitiFact, and Wikipedia for supporting evidence in real time before it makes a decision. This would likely improve performance and reliability on complex or context-dependent misinformation such as satire, propaganda, and subtly manipulated content.

A central limitation of the deep learning models evaluated here, as of deep learning models in general, is that they work largely as black boxes, and an end user cannot easily see why a given decision was made. Future work should add Explainable AI (XAI) methods such as LIME or SHAP for transparency. One promising design is a two-stage pipeline. An LLM, possibly helped by retrieved evidence, first classifies the input, and an explanation layer then produces a human-readable rationale. Such a design is expected to give both strong predictions and interpretability for journalists and fact-checkers, who must act on the model's output.

Beyond LLM integration and explainability, future work should address practical deployment concerns, including:

\begin{itemize}
\item \textbf{Multilingual Detection:} Extending the models beyond English to detect fake news in several languages, especially low-resource languages that existing benchmark datasets underserve.
\item \textbf{Multi-modal Fake News Analysis:} Using images, video, and metadata together with text, since misinformation increasingly spans several media types.
\item \textbf{Federated Learning:} Training models collaboratively across organizations or platforms without centralizing sensitive user or content data.
\item \textbf{API and System Deployment:} Packaging the pipeline as an API or service that can be integrated into social media monitoring tools, browser extensions, or newsroom fact-checking workflows.
\end{itemize}

Pursuing these directions would move the work in this thesis from a controlled, benchmark-driven comparison toward a deployable, explainable, and continually evolving fake news detection system.
